\begin{theorem}\nf{(Предельный переход в равномерно сходящизся рядах)}\label{16_th5}\\
	Если $\ss f_n(x)$ равномерно сходится на метрическом пространстве $E$, $x_0$ - предельная точка $E$, $(\fa n\in\N)\ \lim_{x\ra x_0}f_n(x)=a_n\in\R$, то $\ss a_n$ сходится и $\ss a_n=\lim_{x\ra x_0}^{x\in E}\ss f_n(x)$.
\end{theorem}
\begin{proof}
	Применим (\ref{15_th5}) к $S_n(x)=\sum_{k=1}^nf_k(x)$. Получим требуемое.
\end{proof}
\begin{corollary}
	Если ряд из непрерывных на метрическом пространстве $E$ функций сходится равномерно на $E$, то его сумма непрерывна на $E$.
\end{corollary}
\begin{theorem}\nf{(Признак Дини для рядов)}\label{16_th6}\\
	Если $(\fa n\in\N)\ f_n(x)$ непрерывна и неотрицательна по компактам множества $E$, а $\ss f_n(x)$ сходится к непрерывной на $E$ функции, то $\ss f_n(x)$ сходится равномерно на $E$.
\end{theorem}
\begin{proof}
	Применим (\ref{15_th6}) к $S_n(x)=\sum_{k=1}^nf_k(x)$. Получим требуемое.
\end{proof}
\begin{theorem}\nf{(Интегрирование функциональных последовательностей)}\label{16_th7}\\
	Если $(\fa n\in\N)\ f_n(x)\in\mc R[a,b]$ и $\sq[f_n]$ сходится равномерно на $[a,b]$ к $f(x)$, то:
	\[f(x)\in\mc R[a,b],\ \int_a^bf(x)dx=\lim_{n\ra\i}\int_a^bf_n(x)dx\]
\end{theorem}
\begin{proof}
	\[f_n\rra f\Ra(\fa\eps>0)(\ex N\in\N)(\fa n>N)(\fa x\in[a,b])\ |f_n(x)-f(x)|<\frac\eps{3(b-a)}\]
	Зафиксируем $n$, тогда по критерию интегрируемости:
	\[f_n\in\mc R[a,b]\Ra\ex P:a=x_0<\dots<x_n=b:U(P,f_n)-L(P,f_n)<\frac\eps3\]
	\[U(P,f_n)=\sum_{k=1}^m\sup_{[x_{k-1},x_k]}f_n(x)\cdot\tr x_n\ge\sum_{k=1}^m\l(\sup_{[x_{k-1},x_k]}f(x)-\frac{\eps}{3(b-a)}\r)\cdot\tr x_k=U(P,f)-\frac\eps3\]
	Тогда аналогично:
	\[L(P,f_n)\le L(P,f)+\frac\eps3\]
	Соединив, получаем:
	\[U(P,f)-L(P,f)<\eps\Ra f\in\mc R[a,b]\]
	Теперь покажем значение интеграла.
	\[\l|\int_a^bf(x)dx-\int_a^bf_n(x)dx\r|\le\int_a^b|f(x)-f_n(x)|dx\le\frac\eps3<\eps\]
	Получили требуемое.
\end{proof}
\begin{example}
	\[f_m(x)=\lim_{n\ra\i}(\cos(m!\pi\cdot x))^{2n},\ \lim_{m\ra\i}f_m(x)=D(x)\]
	Заметим, что у $f_m$ лишь конечное число точек разрыва при фиксированном $m$, значит $f_m\in\mc R$ на $\R$, а её предел - нет.
\end{example}
\begin{note}
	Аналогичные утверждения справедливы и для $f_n\in\mc R(\al,[a,b])$ при $f_n\rra f$.
	\[\int_a^bfd\al=\lim_{n\ra\i}\int_a^bf_nd\al\]
	Заметим, если $a_n\rra\al$, то это неверно:
	\[\int_a^nfd\al\neq\lim_{n\ra\i}\int_a^bfd\al_n\]
\end{note}
\begin{theorem}\nf{(Интегрирование функциональных рядов)}\label{16_th7'}
	Если $(\fa n\in\N)\ f_n\in\mc R[a,b],\ \ss f_n(x)$ сходится равномерно на $[a,b]$, то $\ss f_n(x)\in\mc R[a,b]$, причём:
	\[\int_a^b\ss f_n(x)dx=\ss\int_a^bf_n(x)dx\]
\end{theorem}
\begin{theorem}\nf{(Почленное дифференцирование функциональных последовательностей)}\label{16_th8}
	Если:
	\begin{enumerate}
		\item $(\fa n\in\N)\ f_n$ дифференцируема на $(a,b)$.
		\item $(\ex x_0\in(a,b))\ \sq[f_n(x_0)]$ сходится.
		\item $f_n'$ равномерно сходится на $(a,b)$ при $n\ra\i$.
	\end{enumerate}
	то:
	\begin{enumerate}
		\item $f_n\rra f$ на $(a,b)$ при $n\ra\i$.
		\item $f$ дифференцируема на $(a,b)$.
		\item $(\fa x\in(a,b))\ \lim_{n\ra\i}f_n'(x)=f'(x)$.
	\end{enumerate}
\end{theorem}
\begin{note}
	Теорема (\ref{16_th8}) справедлива для любых конечных промежутков, если на концах, где значение включается, говорить об односторонней производной.
\end{note}
\begin{proof}
	Воспользуемся критерием Коши дважды:
	\begin{equation*}
		\begin{split}
			(\fa\eps>0)(\ex N\in\N)(\fa n>N)(\fa p\in\N&)\ |f_{n+p}(x)-f_n(x)|<\frac\eps2\\
			(\fa x\in(a,b)&)\ |f'_{n+p}(x)-f'_n(x)|<\frac\eps{2(b-a)}
		\end{split}
	\end{equation*}
	Возьмём точки $x,t:a<x<t<b$. Тогда по теореме Лагранжа для функции $f_{n+p}(x)-f_n(x)$:
	\begin{equation}\label{16_eq1}
		(f'_{n+p}(\xi)-f'_n(\xi))(x-t)=\big(f_{n+p}(x)-f_n(x)\big)-\big(f_{n+p}(t)-f_n(t)\big)
	\end{equation}
	Применим это к $t=x_0$ при $\fa x\in(a,b)$:
	\begin{multline*}
		\B|f_{n+p}(x)-f_n(x)\B|\le
		\B|\b(f_{n+p}(x)-f_n(x)\b)-\b(f_{n+p}(x_0)-f_n(x_0)\b)\B|+\B|f_{n+p}(x_0)-f_{n+p}(x_0)\B|=\\
		=|f'_{n+p}(\xi_1)-f'_n(\xi_1)|+\B|f_{n+p}(x_0)-f_n(x_0)\B|<\eps
	\end{multline*}
	Получили:
	\[(\fa\eps>0)(\ex N\in\N)(\fa n>N)(\fa p\in\N)(\fa x\in(a,b))\ |f_{n+p}(x)-f_n{x}|<\eps\Ra f_n\rra f\]
	Положим $\ds\phi_n(t):=\frac{f_n(t)-f_n(x)}{t-x},\ \phi(t)=\frac{f(t)-f(x)}{t-x},\ t\in(a,b)\setminus\{x\},\ \lim_{t\ra x}\phi(t)=f'_n(x)$.\\
	Тогда перепишем (\ref{16_eq1}), используя $\phi_n$:
	\[\B|\phi_n(t)-\phi_{n+p}(t)\B|=\B|f'_{n+p}(\xi)-f'_n(\xi)\B|<\frac\eps{2(b-a)}\Ra\phi_n\rra\phi\]
	Тогда:
	\[\lim_{n\ra\i}f'_n(x)=\lim_{t\ra x}\phi(t)=f'(x)\]
	Получили требуемое.
\end{proof}
\begin{example}
	Покажем, что третье требование в теореме (\ref{16_th8}) обязательно.\vspace{0.2cm}\\
	Рассмотрим $\ds f_n(x)=\frac{\sin nx}{\sqrt n}$.
	Зафиксируем $n$, найдем производную:
	\[\l(\frac{\sin nx}{\sqrt n}\r)'=\frac{\cos nx\cdot n\cdot\sqrt n-\frac1{2\sqrt n}\cdot\sin nx}{n}=\cos nx\cdot\sqrt n+\dots\]
	Видим, что в производной есть $\sqrt n$, значит $\nexists\lim_{n\ra\i}f_n'$.
\end{example}
\begin{theorem}\nf{(Почленное дифференцирование функциональных рядов)}\label{16_th8'}
	Если:
	\begin{enumerate}
		\item $f_n$ дифференцируема на $(a,b)$.
		\item $(\ex x_0\in(a,b))\ \ss f_n(x_0)$ сходится.
		\item $\ss f'_n$ равномерно сходится на $(a,b)$.
	\end{enumerate}
	то:
	\begin{enumerate}
		\item $\ss f_n(x)$ сходится равномерно к $S(x)$ на $(a,b)$.
		\item $S(x)$ дифференцируема на $(a,b)$.
		\item $S'(x)=\ss f'_n(x)$.
	\end{enumerate}
\end{theorem}
\begin{theorem}\nf{(Вейерштрасса, пример ван дер Вардена)}\label{16_th9}
	Существует функция, непрерывная на $\R$ и не дифференцируемая ни в одной точке $x\in\R$.
\end{theorem}
\begin{proof}
	\[
	\phi(x)=
	\begin{cases}
		x,\ x\in[0,1]\\
		2-x,\ x\in(1,2]
	\end{cases}
	\ (\fa x\in\R)\ \phi(x+2)=\phi(x)
	\]
	Тогда сама функция:
	\[f(x):=\ss[0]\l(\frac34\r)^n\phi(4^nx)\]
	Заметим:
	\[\l|\l(\frac34\r)\phi(4^nx)\r|\le\l(\frac34\r)^n\]
	Тогда по признаку Вейерштрасса, так как $\ds\ss\l(\frac34\r)^n$ сходится, то и $\ds\ss f_n(x)$ сходится равномерно.
	\[(\fa x\in\R)(\fa m\in\N)(\ex k\in\Z)\ k\le4^mx\le k+1\]
	Зафиксируем $x$, положим $\al_m:=4^{-m}k,\ \beta_m:=4^{-m}(k+1),\ \al_m\le x\le\beta_m,\ \beta_m-\al_m=4^{-m}$.\\
	Хотим показать:
	\[\nexists\lim_{m\ra\i}\frac{f(\beta_m)-f(\al_m)}{\beta_m-\al_m}=\lim_{n\ra\i}g(m)\]
	Рассмотрим результат $\phi(\beta_m)-\phi(\al_m)$ при разных $n$ и $m$:
	\begin{equation*}
		\begin{split}
			n>m&:\ 4^n\beta_m-4^n\al_m=4^{n-m}\cdot(k+1)-4^{n+m}\cdot k\vdots2\Ra\phi(4^n\beta_m)-\phi(4^n\al_m)=0\\
			n=m&:\ \B|\phi(4^m\al_m)-\phi(4^m\al_m)\B|=\B|\phi(k+1)-\phi(k)\B|=1\\
			n<m&:\ \B|\phi(4^n\beta_m)-\phi(4^n\al_m)\B|=\B|4^n\beta_m-4^n\al_m\B|=4^{n-m}
		\end{split}
	\end{equation*}
	Теперь оценим $g(m)$:
	\begin{multline*}
		\l|\frac{f(\beta_m)-f(\al_m)}{\beta_m-\al_m}\r|=4^m\l|\sum_{k=0}^m\l(\frac34\r)^k\B(\phi(4^k\beta_m)-\phi(4^k\al_m)\B)\r|\ge\\
		\ge4^m\l(\l(\frac34\r)^m\cdot1-\sum_{k=0}^{m-1}\l(\frac34\r)^k\cdot4^{k-m}\r)=\\
		=3^m-\sum_{k=0}^{m-1}3^k=3^m-\frac{3^m-1}2\ge\frac12\cdot 3^m   
	\end{multline*}
\end{proof}